\section{Introduction and Boundary}

Finite centered incidence geometry studies finite incidence-defined rows together with the representation or centered-operator channels that make those rows comparable.  The purpose of this paper is to record a bounded grammar for such rows.  It is not a final theory of every FCIG branch, and it does not turn the fingerprint channel into a complete invariant.

The paper is built after public source rows.  The Type-A standard-space paper supplies the first mature theorem row\cite{BabanskyyTypeAFCIG2026}.  The \(S_4\) finite witness paper supplies a small public finite witness for the structured-subset layer\cite{BabanskyyS4Layer2026}.  The Type-A poset-cone and extension-tiler paper supplies the public chamber/tiler context for rows X01--X02\cite{BabanskyyTypeAPoset2026}.  The Latin-square and Sudoku operator papers supply public centered-operator sibling context for rows X05--X06\cite{BabanskyyLatin2026,BabanskyySudokuOperators2026}.  The present paper does not reopen those proofs; it uses them as source rows or related-work rows in a broader grammar.

The organizing principle is that a row is admitted only after it has a finite ambient set, a stated incidence constraint, a representation or operator channel, and a recorded fingerprint package.  Rows without such a finite shadow remain outside the core grammar.

The modular centered row is included as a compact channel exhibit rather than as a new theorem route.  It promotes the reserved X16 row: the same integral centered operator can have one rational augmentation rank and a different rank after reduction modulo a prime.  The key source-locked witness is \(I_7(1,1)\) in characteristic \(2\), where \(2\) does not divide \(7\), the lattice-qualified Smith form is \(\SNFL=[4,4,8]\), and the augmentation rank drops from \(3\) over \(\mathbb Q\) to \(0\) over \(\mathbb F_2\).  The point is that a finite-shadow row may need characteristic, augmentation channel, reduced quotient, and lattice-qualified Smith data, not rational rank alone.  The new dihedral row in this revision is included for that reason: it is a finite vertex action \(I_2(m)\leq S_m\), with explicit incidence rows and replayed standard-block fingerprints.  Its mathematics is elementary cyclic and dihedral linear algebra\cite{BjornerBrenti2005}; its role here is to test whether the grammar still makes useful distinctions in a rank-two reflection shadow.  The length-colored \(G_2\) row is included with the same restraint.  It is not a general Weyl theory; it is a finite root-shadow exhibit showing that the grammar sometimes has to remember source channels, here the short/long partition, rather than only the uncolored ambient permutation set.

The projective \(F_4\) row adds a second, exceptional source-aware test.  It uses the 24-point projective root shadow and the matroid automorphism source of Fried--Gerek--Gordon--Perunicic\cite{FriedGerekGordonPerunicic2007F4Matroid}.  The row is included only to compare two source layers on the same finite set: the geometric projective Weyl layer preserves the long/short split and leaves one centered length residual, while the source matroid automorphism layer adds a non-geometric length-swap coset and cancels that residual.  This is not a theorem about all \(F_4\) rows, Weyl groups, or the 24-cell.

The projective \(H_3\) row is included in the same compact spirit, but with a different channel split.  On the same 15 projective root pairs, the flat-size/matroid channel has automorphism group 120, while the Gram-angle geometric channel has automorphism group 60.  The ordinary standard block is blind to this boundary, and the first higher ordinary Specht separator is \((13,2)\).  The row is used only as a finite source-channel exhibit, not as a theorem about all \(H_3\) shadows or non-crystallographic Coxeter systems.

The projective \(H_4\) row is included only after the selected-atlas source package closes its finite channel checks.  It uses the 60 antipodal pairs in the \(H_4\) root system, viewed as a projective 600-cell root shadow, and compares two pair-color channels: rank-two flat size and squared Gram value.  The flat-size channel has full automorphism group 14400, while the Gram-square channel has full automorphism group 7200.  This is the source-aware lesson used here; it is not a paper about \(H_4\), the 600-cell, or root-system matroid automorphism theory.
The projective \(D_5\) row is included as a control on the pair-channel boundary.  Unlike \(G_2\) or \(F_4\), it has no short/long length split.  The point is that pair-level Gram/flat-size data still has too many automorphisms: \(C_2^{10}:S_5\) of order 122880.  The 40 projective \(A_2\) 3-flat incidences cut this back exactly to the projective Weyl action of order 1920.  This is not a standalone \(D_5\) theorem and not a transfer to \(E_6,E_7,E_8\).

The projective \(E_6\) row is included as a second simply laced control, but with the opposite lesson.  The 36 projective root pairs have pair, flat-size, and circuit channels whose full automorphism groups all agree with the 51840-element projective source action.  The useful grammar datum is therefore not an enlarged automorphism group, but the source-aware decomposition \(\mathbb Q^{\Pi_{36}}=\mathbb Q\mathbf1\oplus U_{20}\oplus U_{15}\): two parabolic rows can have the same size and total standard rank while splitting differently across \(U_{20}\) and \(U_{15}\).  This is not a standalone \(E_6\) theorem and not a Schlaefli, double-six, 27-line, cubic-surface, root-system matroid, classifier, or tiling route.
The projective \(E_7\) row is included as a compact control-plus-separator.  The 63 projective root pairs have Gram-square and rank-two flat-size pair channels whose full automorphism groups already agree with the 1451520-element projective source action.  Nevertheless, embedded row placement is still visible: two \(A_5\) parabolic rows of the same abstract type and order 720 have different centered ranks, \(5=(3,2)\) and \(6=(3,3)\), on \(U_{27}\oplus U_{35}\).  This is not a standalone \(E_7\) theorem, not a reflection-subgroup classification, not a root-system matroid novelty claim, and not an \(E_8\) or Type-E family theorem.
The projective \(E_8\) row is included as the compact endpoint of this E-series comparison.  The 120 projective root pairs have Gram-square and rank-two flat-size pair channels whose full automorphism groups already agree with the 348364800-element projective source action.  Standard simple-parabolic rows of the same abstract type and order fuse under this action, but source-aware subsystem rows can still separate placement: an \(A_7\) parabolic row and an \(A_7\) row inside a coordinate \(D_8\) subsystem have orbit profiles \([8,28,28,56]\) and \([1,28,28,28,35]\), with centered ranks \(3=(1,2)\) and \(4=(1,3)\) on \(U_{35}\oplus U_{84}\).  This is not a standalone \(E_8\) theorem, not a Type-E family theorem, and not a root-subsystem classification.

\paragraph{Contributions.}
The manuscript makes four bounded contributions.
\begin{enumerate}[label=\textup{(\roman*)},leftmargin=*]
  \item It fixes a five-slot finite-shadow grammar \(\mathcal S=(N,X,I,\rho,\PhiFCIG)\), distinguishes family rank from ambient degree, and gives an admission table for the current FCIG source rows.
  \item It uses the elementary standard-block dictionary as an interface lens: position-value balance is exactly the vanishing of the \((N-1,1)\) Specht block, and this lens is applied to both the Type-A row and a signed Type-B/C mirror-union row.
  \item It gives worked finite examples from the Type-A, modular centered, signed Type-B/C, dihedral, length-colored \(G_2\), projective \(F_4\), projective \(H_3\), projective \(H_4\), projective \(D_5\), projective \(E_6\), projective \(E_7\), projective \(E_8\), and centered-operator artifacts where \(\PhiFCIG\) or the source-channel grammar sees data not determined by coarse histograms, rational rank alone, characteristic-zero data alone, uncolored ambient conjugacy alone, or the finite set alone.
  \item It records the guardrails: \(\PhiFCIG\) is descriptive, lattice-qualified where Smith form is used, and never used here as a tiling test or as a substitute for the separate theorem routes.
\end{enumerate}